
Donc il suffit de montrer que l'on a pour tout @@M0@@
@@M1@@
où, comme ci-dessus, @@M2@@, @@M3@@. Or, on a
@@M4@@
pour tout @@M5@@, de sorte que
@@M6@@
Mais (6.2.10) et (6.5.3) donnent
@@M7@@
pour tout @@M8@@, d'où
@@M9@@
avec
@@M10@@
Or, @@M11@@ est une section de @@M12@@, partout égale au coefficient de @@M13@@ dans @@M14@@, donc nulle. On a gagné.

\section*{6. Appendice : Calcul des déterminants des faisceaux localement libres}

\subsection*{6.1.}
Considérons d'abord un topos annelé quelconque @@M15@@. Pour toute section positive @@M16@@ du faisceau constant @@M17@@ et pour tout @@M18@@-module @@M19@@, il est clair comment définir les @@M20@@-modules
